\chapter{Selection and Repetition Statements}
\label{chap:control-flow}

\chapterguide
  {You should be able to construct logical expressions with comparison operators.}
  {use \texttt{if}, \texttt{if...else}, \texttt{if...else if...else}, and \texttt{switch()}; build \texttt{repeat}, \texttt{while}, and \texttt{for} loops; and use \texttt{break} and \texttt{next} in the patterns supplied by the lab.}

\section{Selection statements}

Selection statements determine which block of code executes. The handout introduces them in increasing complexity.

\subsection{Simple \texttt{if}}

\begin{lstlisting}[style=dsr]
x <- 30L

if (is.integer(x)) {
  print("X is an Integer")
}
\end{lstlisting}

The condition is evaluated first. In the supplied example, the message is printed only when \texttt{is.integer(x)} is TRUE.

\subsection{\texttt{if...else}}

\begin{lstlisting}[style=dsr]
x <- c("what", "is", "truth")

if ("Truth" %in% x) {
  print("Truth is found")
} else {
  print("Truth is not found")
}
\end{lstlisting}

This example also reinforces that character comparison is case-sensitive in the supplied code: \texttt{"Truth"} and \texttt{"truth"} are different strings.

\subsection{\texttt{if...else if...else}}

\begin{lstlisting}[style=dsr]
x <- c("what", "is", "truth")

if ("Truth" %in% x) {
  print("Truth is found the first time")
} else if ("truth" %in% x) {
  print("truth is found the second time")
} else {
  print("No truth found")
}
\end{lstlisting}

The conditions are checked in order. This structure is useful when several mutually exclusive outcomes are possible.

\section{\texttt{switch()}}

The handout uses a numeric selector:

\begin{lstlisting}[style=dsr]
x <- switch(
  3,
  "first",
  "second",
  "third",
  "fourth"
)
print(x)
\end{lstlisting}

With selector \texttt{3}, the third listed value is returned.

\section{Repetition statements}

Loops repeat a set of statements. Lab 5a presents three forms.

\subsection{\texttt{repeat} loop}

\begin{lstlisting}[style=dsr]
v <- c("Hello", "loop")
cnt <- 2

repeat {
  print(v)
  cnt <- cnt + 1

  if (cnt > 5) {
    break
  }
}
\end{lstlisting}

The loop itself has no condition in its header. Termination is controlled inside the body using \texttt{break}.

\subsection{\texttt{while} loop}

\begin{lstlisting}[style=dsr]
v <- c("Hello", "while loop")
cnt <- 2

while (cnt < 7) {
  print(v)
  cnt = cnt + 1
}
\end{lstlisting}

The condition is checked before each iteration. The counter update is essential to eventually make the condition false.

\subsection{\texttt{for} loop}

The handout uses \texttt{for} with letters, numbers, and list elements:

\begin{lstlisting}[style=dsr]
v <- LETTERS[1:4]
for (i in v) {
  print(i)
}

for (x in 1:10) {
  print(x)
}

fruits <- list("apple", "banana", "cherry")
for (x in fruits) {
  print(x)
}
\end{lstlisting}

\section{Controlling a loop with \texttt{break} and \texttt{next}}

\texttt{break} exits the current loop:

\begin{lstlisting}[style=dsr]
fruits <- list("apple", "banana", "cherry")

for (x in fruits) {
  if (x == "cherry") {
    break
  }
  print(x)
}
\end{lstlisting}

\texttt{next} skips the remainder of the current iteration and continues with the next element:

\begin{lstlisting}[style=dsr]
for (x in fruits) {
  if (x == "banana") {
    next
  }
  print(x)
}
\end{lstlisting}

\section{Conditions inside loops}

The dice example combines iteration, a comparison, and an \texttt{if...else} block:

\begin{lstlisting}[style=dsr]
dice <- 1:6

for (x in dice) {
  if (x == 6) {
    print(paste("The dice number is", x, "Yahtzee!"))
  } else {
    print(paste("The dice number is", x, "Not Yahtzee"))
  }
}
\end{lstlisting}

This is a useful model for many practical problems: iterate through values, test each value, and produce an output based on the condition.

\section{Nested loops}

The final Lab 5a example places one \texttt{for} loop inside another:

\begin{lstlisting}[style=dsr]
adj <- list("red", "big", "tasty")
fruits <- list("apple", "banana", "cherry")

for (x in adj) {
  for (y in fruits) {
    print(paste(x, y))
  }
}
\end{lstlisting}

Each adjective is paired with each fruit. The inner loop completes its cycle for every value of the outer loop.

\begin{quickcheck}
\begin{enumerate}
  \item What must be TRUE for the body of a simple \texttt{if} statement to run?
  \item Why does the first branch of the handout's \texttt{"Truth"} example fail while the second succeeds?
  \item How does a \texttt{repeat} loop terminate in the supplied pattern?
  \item What is the difference between \texttt{break} and \texttt{next} in the fruit examples?
  \item How many adjective-fruit combinations are generated by three adjectives and three fruits?
\end{enumerate}
\end{quickcheck}

\begin{chapterrecap}
\begin{itemize}
  \item Selection statements choose among code paths using logical conditions.
  \item \texttt{switch()} can select a result by position in the supplied example.
  \item \texttt{repeat} relies on an internal \texttt{break}; \texttt{while} checks a condition before each iteration; \texttt{for} iterates over supplied values.
  \item \texttt{break} stops a loop and \texttt{next} skips one iteration.
  \item Loops and selection statements can be nested or combined to solve multi-step tasks.
\end{itemize}
\end{chapterrecap}
